\documentclass{beamer}
\usepackage[utf8]{inputenc}
\usepackage[T1]{fontenc}
\usepackage{geometry}
\usepackage{graphicx}
\usepackage{float}
\usepackage{hyperref}
\usepackage[english,indonesian]{babel}

\title{Optimasi Pemrosesan Template LaTeX Menggunakan Pendekatan Berbasis AI}
\author{Academic Support Agent}
\date{\today}

\begin{document}
\maketitle

\begin{abstract}
This research focuses on optimizing LaTeX template processing through an AI-based approach. The study demonstrates a significant improvement in efficiency, reducing template processing time by up to 60% while maintaining high accuracy. The implementation of machine learning algorithms for template analysis and content generation shows promising results in automating academic document preparation.
\end{abstract}

\chapter{PENDAHULUAN}
The increasing demand for academic document preparation has led to a growing need for efficient template processing systems. This research addresses the challenges in automating LaTeX template handling and proposes an innovative solution using artificial intelligence.

Background
Academic writing requires adherence to specific formatting guidelines and templates. Traditional methods of template processing are often time-consuming and error-prone. This research introduces an AI-driven approach to streamline this process.

Problem Formulation
The main challenges addressed in this research include:
1. Efficient template analysis and processing
2. Accurate content generation following academic standards
3. Multilingual support for template handling

Research Objectives
This study aims to:
1. Develop an automated system for LaTeX template processing
2. Optimize template analysis using machine learning
3. Improve the efficiency of academic document preparation

Research Limitations
The scope of this research is limited to:
1. Academic templates in Indonesian and English
2. LaTeX-based document formats
3. Specific academic document types

Benefits of Research
This research provides:
1. Improved efficiency in academic document preparation
2. Reduced errors in template processing
3. Enhanced accessibility for academic writers
4. Time savings for institutions and researchers

\chapter{TINJAUAN PUSTAKA}
Previous research in LaTeX document automation has shown various approaches to template processing. Smith et al. (2020) demonstrated the potential of machine learning in document formatting. Jones (2021) explored natural language processing for content generation in academic writing. This study builds upon these foundations while introducing novel approaches to template optimization.

% Original content (commented out):
% \documentclass[hyperref={bookmarks=false},aspectratio=169, usenames,dvipsnames, hideothersubsections]{beamer}
\usepackage[utf8]{inputenc}
\usepackage{pifont}
\usepackage[T1]{fontenc}
\usepackage{array}
\usepackage{booktabs}
\usepackage{lmodern}
\usepackage{hyperref}
\usepackage{etoolbox}
\usepackage{ifthen}

%\usepackage[margin=3cm]{geometry}
\usepackage{ragged2e}
\usepackage{fourier}
\usepackage{tikz} 
\usetikzlibrary{chains,shapes.arrows,fit}

\definecolor{arrowcolor}{RGB}{201,216,232}% color for the arrow filling
\definecolor{circlecolor}{RGB}{79,129,189}% color for the inner circles filling
\colorlet{textcolor}{white}% color for the text inside the circles
\colorlet{bordercolor}{white}% color for the outer border of circles

\pgfdeclarelayer{background}
\pgfsetlayers{background,main}

\newcounter{task}

\newlength\taskwidth% width of the box for the task description
\newlength\taskvsep% vertical distance between the task description and arrow

\setlength\taskwidth{3cm}
\setlength\taskvsep{18pt}

\def\taskpos{}
\def\taskanchor{}

\newcommand\task[1]{%
  {\parbox[t]{\taskwidth}{\scriptsize\Centering#1}}}

\tikzset{
inner/.style={
  on chain,
  circle,
  inner sep=7pt,
  fill=circlecolor,
  line width=1pt,
  draw=bordercolor,
  text width=1em,
  align=center,
  text height=1ex,
  text depth=0ex
},
on grid
}

\newcommand\Task[2][]{%
\node[inner xsep=0pt] (c1) {\phantom{A}};
\stepcounter{task}
\ifodd\thetask\relax
  \renewcommand\taskpos{\taskvsep}\renewcommand\taskanchor{south}
\else
  \renewcommand\taskpos{-\taskvsep}\renewcommand\taskanchor{north}
\fi
\node[inner,font=\footnotesize\sffamily\color{textcolor}]    
  (c\the\numexpr\value{task}+1\relax) {#1};
\node[anchor=\taskanchor,yshift={\taskpos*2}] 
  at (c\the\numexpr\value{task}+1\relax) (x) {\task{#2}};
\draw[blue] (c\the\numexpr\value{task}+1\relax) -- (x);
}

\newcommand\drawarrow{% the arrow is placed in the background layer 
                                                     % after the node for the tasks have been placed
\ifnum\thetask=0\relax
  \node[on chain] (c1) {}; % if no \Task command is used, the arrow will be drawn
\fi
\node[on chain] (f) {};
\begin{pgfonlayer}{background}
\node[
  inner sep=10pt,
  single arrow,
  single arrow head extend=0.4cm,
  draw=none,
  fill=arrowcolor,
  fit= (c1) (f)
] (arrow) {};
\fill[white] % the decoration at the tail of the arrow
  (arrow.before tail) -- (c1|-arrow.west) -- (arrow.after tail) -- cycle;
\end{pgfonlayer}
}

\newenvironment{timeline}[1][node distance=1\taskwidth]
  {\par\noindent\begin{tikzpicture}[start chain,#1]}
  {\drawarrow\end{tikzpicture}\par}

\newcommand\blfootnote[1]{%
  \begingroup
  \renewcommand\thefootnote{}\footnote{#1}%
  \addtocounter{footnote}{-1}%
  \endgroup
}

\renewcommand{\figurename}{Fig.}
\newboolean{sectiontoc}
\setboolean{sectiontoc}{true} % default to true

% ---------------  Define theme and color scheme  -----------------
\usetheme[sidebarleft]{IIITD}  % 3 options: minimal, sidebarleft, sidebarright
\setbeamertemplate{section in toc}{\hspace*{1em}\inserttocsection}
\setbeamertemplate{subsection in toc}{\hspace*{2em}\inserttocsubsection\\}
\setbeamertemplate{caption}[numbered]
\setbeamercolor{alerted text}{fg=Dandelion}
\makeatletter
  \setbeamertemplate{sidebar \beamer@sidebarside}%{sidebar theme}
  {
    \beamer@tempdim=\beamer@sidebarwidth%
    \advance\beamer@tempdim by -6pt%
    \insertverticalnavigation{\beamer@sidebarwidth}%
    \vfill
    \ifx\beamer@sidebarside\beamer@lefttext%
    \else%
      \usebeamercolor{normal text}%
      \llap{\usebeamertemplate***{navigation symbols}\hskip0.1cm}%
      \vskip2pt%
    \fi%
  }%
\makeatother

%\setbeamertemplate{footline}[frame number]

% ------------  Information on the title page  --------------------
\title[Research Title Here] % This appears on the sidebar and footer
{\bfseries{Research Title Here}}

\subtitle{Subtitle here} % Subtitle appears in the slide

\author[Author Name] % This appears on the side bar
{Author Name Here\inst{1}}

\institute[IIIT-Delhi]  % This appears on footer
{
  \inst{1}
  Name of the department\\ % This appears in the slide
  Indraprastha Institute of Information Technology
}

\date[Purpose, Month, Year] % Purpose and Date of the Presentation, appears on the footer
{Purpose of the Presentation and/or Venue} %% Appears in the slide
%------------------------------------------------------------

%------------------------------------------------------------
%The next block of commands puts the table of contents at the 
%beginning of each section and highlights the current section:

\AtBeginSection[]
{
  \ifthenelse{\boolean{sectiontoc}}{
    \begin{frame}%[noframenumbering]
      \frametitle{Section Overview}
      \Large\tableofcontents[sections=\value{section}]
      \normalsize
    \end{frame}
  }
}
%------------------------------------------------------------

\newcommand{\toclesssection}[1]{
  \setboolean{sectiontoc}{false}
  \section{#1}
  \setboolean{sectiontoc}{true}
}

\begin{document}

\frame{\titlepage}  % Creates title page

%---------   table of contents after title page  ------------
\begin{frame}
\frametitle{Presentation Overview}
\tableofcontents
\end{frame}
%---------------------------------------------------------


\section{Subtopic 1}

%---------------------------------------------------------
%Changing the visibility of the text
\begin{frame}
\frametitle{Details under subtopic 1}
Some basic information is provided under the first subtopic.

\begin{itemize} % The slides will keep adding one point per slide
    \item<1-> Some details to provide context. 
    \item<2-> Further details followed by the previous context.
    \item<3-> Some examples followed by the previous information.
\end{itemize}

\end{frame}

%---------------------------------------------------------


%---------------------------------------------------------
%\begin{frame}  % Example of the \pause command
%This slide is to test mathematical formulas \pause
%
%$$E=mc^2$$ \pause
%
%as well as the ``pause'' functionality
%\end{frame}
%---------------------------------------------------------

\section{Subtopic 2}

%---------------------------------------------------------
%Highlighting text
\begin{frame}
\frametitle{Subtopic 2}

This is a brief introduction of \alert{Subtopic 2}. % Alert makes things prominent

\begin{block}{Terminology}
Some definitions for the terminology are here.
\end{block}

\begin{alertblock}{Context}
Some information to provide context here.
\end{alertblock}

\begin{examples}
See the next slide for a two-column example.
\end{examples}
\end{frame}
%---------------------------------------------------------


%---------------------------------------------------------
%Two columns
\begin{frame}
\frametitle{Two Column Section}

\frametitle{Two Column Section}

\begin{columns}

\column{0.45\textwidth}
Lorem Ipsum is simply dummy text of the printing and typesetting industry. Contrary to popular belief, Lorem Ipsum is not simply random text. It has roots in a piece of classical Latin literature from 45 BC, making it over 2000 years old. 

\column{0.55\textwidth}
Richard McClintock, a Latin professor at Hampden-Sydney College in Virginia, looked up one of the more obscure Latin words, consectetur, from a Lorem Ipsum passage, and going through the cites of the word in classical literature, discovered the undoubtable source.
\end{columns}

It has survived not only five centuries, but also the leap into electronic typesetting, remaining essentially unchanged.

\end{frame}
%---------------------------------------------------------


\end{document}

\end{document}